% GENERATED FILE. Edit canonical lessons, then run npm run generate:books.
% canonical-source-hash: fnv1a64:31be16dfaf3d7d5f
% canonical-lessons: UR-C547-tasavvur, UR-C547-niyyat, UR-C547-muhim, UR-C547-peshkash, UR-C547-sanvarna-intr, UR-R547-review-541-547, UR-R547-review-final-547

\chapter{Looking Ahead}
\label{ch:ur-looking-ahead}
\hlchaptermodality{547}

\begin{chapteropening}
\textbf{By the end of this chapter:} \emph{I can say what might happen in the future if things go well.}

Complete the last lesson of chapter 547: I can say what might happen in the future if things go well.
\end{chapteropening}

\section[\texorpdfstring{\ur{تصور} (tasavvur)}{تصور (tasavvur)}]{\ur{تصور} (tasavvur) — imagination; a picture in the mind}

\label{lesson:UR-C547-tasavvur}



\begin{quote}
Before the new one: say the Urdu for to grapple with (a hard task), then the Urdu for a yardstick.
\end{quote}

\subsection*{You'll want to know: \ur{تصور}}
\textbf{\ur{تصور}} — \emph{tasavvur} — \textquotedblleft{}imagination; a picture in the mind\textquotedblright{}.

\begin{quote}
\ur{اگر} \ur{تصور} \ur{نہ} \ur{ہوتا} \ur{تو} \ur{کوئی} \ur{نئی} \ur{چیز} \ur{نہ} \ur{بنتی۔}
\end{quote}

\begin{quote}
\emph{agar tasavvur na hotā to koī naī chīz na bantī.}
\end{quote}

\begin{quote}
Without imagination nothing new would ever be made.
\end{quote}

Masculine.

\begin{cousinweb}
\textbf{\ur{تصور}} (\emph{tasavvur}) is usually traced, through Persian, to Arabic \emph{taṣawwur}, \textquotedblleft{}forming an image in the mind\textquotedblright{}.
\end{cousinweb}

\begin{grammarlens}[title={What you've built}]
One more word for saying what happens if.
\end{grammarlens}

\subsection*{Your turn}
\begin{itemize}
\raggedright
  \item \emph{Say it:} \emph{tasavvur}
  \item \emph{Say it:} \emph{tasavvur}, once more
  \item \emph{Say it:} \emph{agar tasavvur na hotā to koī naī chīz na bantī.} — Without imagination nothing new would ever be made.
  \item \emph{Recall:} say the Urdu for to grapple with (a hard task), then the Urdu for a yardstick, then say \emph{tasavvur} again
\end{itemize}

\begin{tcolorbox}[breakable,colback=teal!4,colframe=teal!35!black,title={Before you move on}]
What does \ur{تصور} mean? (\textquotedblleft{}Imagination; a picture in the mind\textquotedblright{}.) Say it once more.
\end{tcolorbox}
\section[\texorpdfstring{\ur{نیت} (niyyat)}{نیت (niyyat)}]{\ur{نیت} (niyyat) — an intention (in one's heart)}

\label{lesson:UR-C547-niyyat}



\begin{quote}
Before the new one: say the Urdu for a yardstick, then the Urdu for imagination.
\end{quote}

\subsection*{You'll want to know: \ur{نیت}}
\textbf{\ur{نیت}} — \emph{niyyat} — \textquotedblleft{}an intention (in one's heart)\textquotedblright{}.

\begin{quote}
\ur{اگر} \ur{نیت} \ur{اچھی} \ur{ہے} \ur{تو} \ur{راستہ} \ur{مل} \ur{جائے} \ur{گا۔}
\end{quote}

\begin{quote}
\emph{agar niyyat achchhī hai to rāstā mil jāe gā.}
\end{quote}

\begin{quote}
If your heart is in the right place, you'll find a way.
\end{quote}

Feminine.

\begin{cousinweb}
\textbf{\ur{نیت}} (\emph{niyyat}) is usually traced, through Persian, to Arabic \emph{niyya}, \textquotedblleft{}an intention\textquotedblright{}.
\end{cousinweb}

\begin{grammarlens}[title={What you've built}]
One more word for saying what happens if.
\end{grammarlens}

\subsection*{Your turn}
\begin{itemize}
\raggedright
  \item \emph{Say it:} \emph{niyyat}
  \item \emph{Say it:} \emph{niyyat}, once more
  \item \emph{Say it:} \emph{agar niyyat achchhī hai to rāstā mil jāe gā.} — If your heart is in the right place, you'll find a way.
  \item \emph{Recall:} say the Urdu for a yardstick, then the Urdu for imagination, then say \emph{niyyat} again
\end{itemize}

\begin{tcolorbox}[breakable,colback=teal!4,colframe=teal!35!black,title={Before you move on}]
What does \ur{نیت} mean? (\textquotedblleft{}An intention (in one's heart)\textquotedblright{}.) Say it once more.
\end{tcolorbox}
\section[\texorpdfstring{\ur{مہم} (muhim)}{مہم (muhim)}]{\ur{مہم} (muhim) — a campaign; a venture}

\label{lesson:UR-C547-muhim}



\begin{quote}
Before the new one: say the Urdu for imagination, then the Urdu for an intention (in one's heart).
\end{quote}

\subsection*{You'll want to know: \ur{مہم}}
\textbf{\ur{مہم}} — \emph{muhim} — \textquotedblleft{}a campaign; a venture\textquotedblright{}.

\begin{quote}
\ur{اگر} \ur{مہم} \ur{ختم} \ur{ہوئی} \ur{تو} \ur{سب} \ur{گھر} \ur{لوٹیں} \ur{گے۔}
\end{quote}

\begin{quote}
\emph{agar muhim khatm huī to sab ghar lauṭeṅ ge.}
\end{quote}

\begin{quote}
If the campaign ends, everyone will go home.
\end{quote}

Feminine.

\begin{cousinweb}
\textbf{\ur{مہم}} (\emph{muhim}) is usually traced, through Persian, to Arabic \emph{muhimm}, \textquotedblleft{}an important matter\textquotedblright{}; the sense \textquotedblleft{}a campaign\textquotedblright{} came later, and the sources we use do not show how it arose.
\end{cousinweb}

\begin{grammarlens}[title={What you've built}]
One more word for saying what happens if.
\end{grammarlens}

\subsection*{Your turn}
\begin{itemize}
\raggedright
  \item \emph{Say it:} \emph{muhim}
  \item \emph{Say it:} \emph{muhim}, once more
  \item \emph{Say it:} \emph{agar muhim khatm huī to sab ghar lauṭeṅ ge.} — If the campaign ends, everyone will go home.
  \item \emph{Recall:} say the Urdu for imagination, then the Urdu for an intention (in one's heart), then say \emph{muhim} again
\end{itemize}

\begin{tcolorbox}[breakable,colback=teal!4,colframe=teal!35!black,title={Before you move on}]
What does \ur{مہم} mean? (\textquotedblleft{}A campaign; a venture\textquotedblright{}.) Say it once more.
\end{tcolorbox}
\section[\texorpdfstring{\ur{پیشکش} (peshkash)}{پیشکش (peshkash)}]{\ur{پیشکش} (peshkash) — an offer}

\label{lesson:UR-C547-peshkash}



\begin{quote}
Before the new one: say the Urdu for an intention (in one's heart), then the Urdu for a campaign.
\end{quote}

\subsection*{You'll want to know: \ur{پیشکش}}
\textbf{\ur{پیشکش}} — \emph{peshkash} — \textquotedblleft{}an offer\textquotedblright{}.

\begin{quote}
\ur{اگر} \ur{یہ} \ur{پیشکش} \ur{اچھی} \ur{ہوئی} \ur{تو} \ur{میں} \ur{قبول} \ur{کروں} \ur{گا۔}
\end{quote}

\begin{quote}
\emph{agar yih peshkash achchhī huī to maiṅ qabūl karūṅ gā.}
\end{quote}

\begin{quote}
If the offer is good, I will accept it.
\end{quote}

Feminine.

\begin{cousinweb}
\textbf{\ur{پیشکش}} (\emph{peshkash}) is usually traced to Persian \emph{pēshkash}, \textquotedblleft{}something offered\textquotedblright{}, built from \emph{pēsh}, \textquotedblleft{}before\textquotedblright{}, and \emph{kash-}, \textquotedblleft{}to draw\textquotedblright{}.
\end{cousinweb}

\begin{grammarlens}[title={What you've built}]
One more word for saying what happens if.
\end{grammarlens}

\subsection*{Your turn}
\begin{itemize}
\raggedright
  \item \emph{Say it:} \emph{peshkash}
  \item \emph{Say it:} \emph{peshkash}, once more
  \item \emph{Say it:} \emph{agar yih peshkash achchhī huī to maiṅ qabūl karūṅ gā.} — If the offer is good, I will accept it.
  \item \emph{Recall:} say the Urdu for an intention (in one's heart), then the Urdu for a campaign, then say \emph{peshkash} again
\end{itemize}

\begin{tcolorbox}[breakable,colback=teal!4,colframe=teal!35!black,title={Before you move on}]
What does \ur{پیشکش} mean? (\textquotedblleft{}An offer\textquotedblright{}.) Say it once more.
\end{tcolorbox}
\section[\texorpdfstring{\ur{سنورنا} (sanvarnā)}{سنورنا (sanvarnā)}]{\ur{سنورنا} (sanvarnā) — to take shape nicely; to be put right}

\label{lesson:UR-C547-sanvarna-intr}



\begin{quote}
Before the new one: say the Urdu for a campaign, then the Urdu for an offer.
\end{quote}

\subsection*{You'll want to know: \ur{سنورنا}}
\textbf{\ur{سنورنا}} — \emph{sanvarnā} — \textquotedblleft{}to take shape nicely; to be put right\textquotedblright{}.

\begin{quote}
\ur{اگر} \ur{بچوں} \ur{نے} \ur{اچھا} \ur{سیکھا} \ur{تو} \ur{کل} \ur{سنور} \ur{جائے} \ur{گا۔}
\end{quote}

\begin{quote}
\emph{agar bachchoṅ ne achchhā sīkhā to kal sanvar jāe gā.}
\end{quote}

\begin{quote}
If the children learn well, tomorrow will turn out well.
\end{quote}

Past (it …): \ur{وہ} \ur{کمرہ} \ur{سنور} \ur{گیا}, \emph{voh kamrā sanvar gayā}.

\begin{cousinweb}
\textbf{\ur{سنورنا}} (\emph{sanvarnā}) is inherited: CDIAL lists the Hindi cousins \emph{sãwārṇā}, \textquotedblleft{}to arrange, to deck\textquotedblright{}, and its intransitive \emph{sãwarnā}, \textquotedblleft{}to be arranged\textquotedblright{}, under Sanskrit \emph{saṃvārayati}, \textquotedblleft{}puts in order\textquotedblright{} (CDIAL 13021).
\end{cousinweb}

\begin{grammarlens}[title={What you've built}]
One more word for saying what happens if.
\end{grammarlens}

\subsection*{Your turn}
\begin{itemize}
\raggedright
  \item \emph{Say it:} \emph{sanvarnā}
  \item \emph{Say it:} \emph{sanvarnā}, once more
  \item \emph{Say it:} \emph{agar bachchoṅ ne achchhā sīkhā to kal sanvar jāe gā.} — If the children learn well, tomorrow will turn out well.
  \item \emph{Recall:} say the Urdu for a campaign, then the Urdu for an offer, then say \emph{sanvarnā} again
\end{itemize}

\begin{tcolorbox}[breakable,colback=teal!4,colframe=teal!35!black,title={Before you move on}]
What does \ur{سنورنا} mean? (\textquotedblleft{}To take shape nicely; to be put right\textquotedblright{}.) Say it once more.
\end{tcolorbox}
\section[(dialogue)]{Looking back: from In That Case to Looking Ahead}

\label{lesson:UR-R547-review-541-547}



\begin{quote}
Say the Urdu for an angle and circumstances.
\end{quote}

\subsection*{Your turn}
Say these aloud:

\begin{itemize}
\raggedright
  \item an angle, circumstances, conditional, a point (that things come to), supposing (for argument's sake)
  \item in any case, well (moving on), in particular, a point (in an argument), a summary
  \item what one expects, a prediction, to see through, to put a value on, correct
  \item a beginning, the start, the close (of an event), a series, order (a set sequence)
  \item temporary, permanent, final, a stage, durable
  \item toil, to lend a hand, ease, to grapple with (a hard task), a yardstick
  \item imagination, an intention (in one's heart), a campaign, an offer, to take shape nicely
\end{itemize}

\begin{tcolorbox}[breakable,colback=teal!4,colframe=teal!35!black,title={Before you move on}]
An angle? (\textbf{\ur{زاویہ}}, \emph{zāviyā}.) Circumstances? (\textbf{\ur{حالات}}, \emph{hālāt}.) Conditional? (\textbf{\ur{مشروط}}, \emph{mashrūt}.) A point (that things come to)? (\textbf{\ur{نوبت}}, \emph{naubat}.) Supposing (for argument's sake)? (\textbf{\ur{بالفرض}}, \emph{bilfarz}.) In any case? (\textbf{\ur{بہرحال}}, \emph{baharhāl}.) Well (moving on)? (\textbf{\ur{خیر}}, \emph{khair}.) In particular? (\textbf{\ur{بالخصوص}}, \emph{bilkhusūs}.) A point (in an argument)? (\textbf{\ur{نکتہ}}, \emph{nuktā}.) A summary? (\textbf{\ur{خلاصہ}}, \emph{khulāsā}.) What one expects? (\textbf{\ur{توقع}}, \emph{tavaqqo}.) A prediction? (\textbf{\ur{پیش} \ur{گوئی}}, \emph{pesh goī}.) To see through? (\textbf{\ur{تاڑنا}}, \emph{tāṛnā}.) To put a value on? (\textbf{\ur{آنکنا}}, \emph{āṅknā}.) Correct? (\textbf{\ur{درست}}, \emph{durust}.) A beginning? (\textbf{\ur{آغاز}}, \emph{āghāz}.) The start? (\textbf{\ur{شروع}}, \emph{shurū}.) The close (of an event)? (\textbf{\ur{اختتام}}, \emph{ikhtitām}.) A series? (\textbf{\ur{سلسلہ}}, \emph{silsilā}.) Order (a set sequence)? (\textbf{\ur{ترتیب}}, \emph{tartīb}.) Temporary? (\textbf{\ur{عارضی}}, \emph{ārzī}.) Permanent? (\textbf{\ur{مستقل}}, \emph{mustaqil}.) Final? (\textbf{\ur{حتمی}}, \emph{hatmī}.) A stage? (\textbf{\ur{مرحلہ}}, \emph{marhalā}.) Durable? (\textbf{\ur{پائیدار}}, \emph{pāedār}.) Toil? (\textbf{\ur{مشقت}}, \emph{mashaqqat}.) To lend a hand? (\textbf{\ur{ہاتھ} \ur{بٹانا}}, \emph{hāth baṭānā}.) Ease? (\textbf{\ur{آسانی}}, \emph{āsānī}.) To grapple with (a hard task)? (\textbf{\ur{جوجھنا}}, \emph{jūjhnā}.) A yardstick? (\textbf{\ur{پیمانہ}}, \emph{paimānā}.) Imagination? (\textbf{\ur{تصور}}, \emph{tasavvur}.) An intention (in one's heart)? (\textbf{\ur{نیت}}, \emph{niyyat}.) A campaign? (\textbf{\ur{مہم}}, \emph{muhim}.) An offer? (\textbf{\ur{پیشکش}}, \emph{peshkash}.) To take shape nicely? (\textbf{\ur{سنورنا}}, \emph{sanvarnā}.)
\end{tcolorbox}
\section[(dialogue)]{Looking Ahead — once more}

\label{lesson:UR-R547-review-final-547}



\begin{quote}
Say the Urdu for imagination and an intention (in one's heart).
\end{quote}

\subsection*{Your turn}
\begin{itemize}
\raggedright
  \item \emph{Say it:} imagination, an intention (in one's heart), a campaign, an offer, to take shape nicely
\end{itemize}

\begin{tcolorbox}[breakable,colback=teal!4,colframe=teal!35!black,title={Before you move on}]
Imagination? (\textbf{\ur{تصور}}, \emph{tasavvur}.) An intention (in one's heart)? (\textbf{\ur{نیت}}, \emph{niyyat}.) A campaign? (\textbf{\ur{مہم}}, \emph{muhim}.) An offer? (\textbf{\ur{پیشکش}}, \emph{peshkash}.) To take shape nicely? (\textbf{\ur{سنورنا}}, \emph{sanvarnā}.)
\end{tcolorbox}
